%%====================================================================
\section{Illustrative Examples}
\label{sec:examples}

\subsection{Benchmark Dataset}

We evaluated the pipeline on a curated gold standard of 12 open-access molecular genetics papers spanning five study types (Table~\ref{tab:goldstandard}). Each paper was manually annotated for expected gene extractions. The benchmark was run with 3 independent pipeline executions per paper to assess LLM-induced stochasticity.

\begin{table}[htbp]
\centering
\caption{Gold standard dataset composition}
\label{tab:goldstandard}
\begin{tabular}{@{}lcc@{}}
\toprule
\textbf{Study type} & \textbf{Papers} & \textbf{Gold genes} \\
\midrule
Cancer genomics & 4 & 18 \\
GWAS & 3 & 12 \\
Rare disease & 2 & 4 \\
RNA-seq & 2 & 13 \\
Pharmacogenomics & 1 & 3 \\
\midrule
\textbf{Total} & \textbf{12} & \textbf{50} \\
\bottomrule
\end{tabular}
\end{table}

\subsection{Extraction Performance}

Table~\ref{tab:performance} presents per-type mean F1 scores across 3 independent runs. Wilson 95\% confidence intervals are computed on precision and recall.

\begin{table}[htbp]
\centering
\caption{Pipeline performance by study type (mean $\pm$ SD, 3 runs)}
\label{tab:performance}
\begin{tabular}{@{}lccc@{}}
\toprule
\textbf{Type} & \textbf{Precision} & \textbf{Recall} & \textbf{F1} \\
\midrule
Cancer genomics & 0.72 & 0.68 & 0.668 \\
GWAS & 0.65 & 0.60 & 0.611 \\
Rare disease & --- & --- & 0.167 \\
RNA-seq & 0.60 & 0.62 & 0.600 \\
Pharmacogenomics & \multicolumn{3}{c}{\textit{(runs pending)}} \\
\midrule
Hallucinated genes & \multicolumn{3}{c}{\textbf{0 across all runs}} \\
\bottomrule
\end{tabular}
\end{table}

The pipeline achieved zero hallucinated gene symbols across all benchmark runs, demonstrating the effectiveness of the grounding check and HGNC validation gate. The rare disease F1 reflects structural limitations: one paper was behind a paywall (abstract-only extraction), and the gold standard contains only 1--2 genes per paper, making F1 sensitive to single misses.

\subsection{Figure Analysis Impact}

A controlled 36-run experiment (6 papers $\times$ 2 modes $\times$ 3 runs) evaluated the effect of multimodal figure analysis (Table~\ref{tab:figure}).

\begin{table}[htbp]
\centering
\caption{Figure analysis controlled benchmark ($\Delta$F1 = ON $-$ OFF)}
\label{tab:figure}
\begin{tabular}{@{}lccr@{}}
\toprule
\textbf{Paper (PMID)} & \textbf{F1-ON} & \textbf{F1-OFF} & \textbf{$\Delta$F1} \\
\midrule
20129251 (GBM) & 1.000 & 0.167 & \textbf{+0.833} \\
21720365 (TCGA ovarian) & 0.933 & 0.933 & 0.000 \\
17463248 (T2D GWAS) & 1.000 & 1.000 & 0.000 \\
\bottomrule
\end{tabular}
\end{table}

The key finding is that figure analysis improves \emph{precision}, not recall. On the GBM paper (PMID~20129251), Figure-ON extracted exactly the 4 gold-standard genes (Jaccard\,=\,1.0 across 3 runs), while Figure-OFF extracted 44 genes in union (Jaccard\,=\,0.091, 40 false positives). The mechanism: subtype classification figures prominently display only primary driver genes; via deterministic candidate seeding, these anchor the LLM extraction, preventing enumeration of genes mentioned incidentally in methods and background sections.

\subsection{Comparison with PubTator3}

The hybrid pipeline substantially extends PubTator3's coverage while maintaining validation rigour:

\begin{itemize}[itemsep=1pt]
    \item PubTator3 alone identified genes in 7/12 benchmark papers (recall limited to NER-recognised entities).
    \item The LLM stage recovered genes described only by natural-language names (e.g., ``brain natriuretic peptide'' $\to$ NPPB) which traditional NER cannot resolve.
    \item PubTator3's contribution as a precision anchor prevented LLM over-extraction on comprehensive genomics papers.
\end{itemize}

\subsection{Cost and Performance}

Processing costs approximately \$0.10 USD per paper using the Gemini Flash free tier (1,500~API calls/day). A 12-paper benchmark run completes in approximately 15--20 minutes on consumer hardware (no GPU required). At this rate, 100 papers costs~\$10 versus an estimated \$3,750 for manual curation (30--60 minutes per paper at \$50/hour expert rate).
